\scanpage{100}
% archive foliation: 99
II) Supposons $\mathcal C_0$ fini. Pour un schéma de base $T$,

\deleted{Appelons revêtement spécial $\mathcal C_0$-épinglé de $\mathbb P^1_T$, un rev. fini \inserted{$X\xrightarrow{f}\mathbb P^1_T$} de $\mathbb P^1_T$, étale sur $U_T=\mathbb P^1_T-\{0,1,\infty\}$, avec \deleted{\unsure{une}}\inserted{ramification} modérée fibre par fibre, et \deleted{des \illegible{1}}\inserted{isomorphismes}}
\[
\deleted{f^{-1}(0_T)_{\unsure{\mathrm{red}}}\simeq}
\]
\[
\begin{aligned}
f^{-1}(0_T)_{\mathrm{réd}}&\simeq\deleted{\unsure{T}}\ S_{T\,\deleted{\unsure{\mathrm{ét}}}\,\mathrm{réd}}\\
f^{-1}(1_T)_{\mathrm{réd}}&\simeq A_{T\,\deleted{\unsure{\mathrm{ét}}}\,\mathrm{réd}}\\
f^{-1}(\infty_T)_{\mathrm{réd}}&\simeq F_{T\,\deleted{\unsure{\mathrm{réd}}}}
\end{aligned}
\]
\deleted{qui fibre par fibre sont \unsure{des isom.} \unsure{précisés} par}
\inserted{``compatibles avec les degrés de ramification'' $(p,q,r,\ \text{en bas})$}
\margin{\deleted{\unsure{fibre par fibre}\ \unsure{évidente}\ \unsure{isomorphismes}\ \unsure{aux fibres}\ \unsure{uniforme}}}

\deleted{\unsure{De plus}} On voit que \deleted{cette} catégorie est rigide.\margin{vérifier}
\deleted{Pour $\mathcal C_0$ fini,} Soit \deleted{$\unsure{F}$} $P_{\mathcal C_0}(S)$ l'ens des classes d'iso. de cette catégorie. On voit que pour $S$ variable, on a un foncteur contravariant en $S$, qui est représentable. On trouve que le schéma $M_{\mathcal C_0}$ qui le représente est quasi fini, \deleted{et} étale et séparé sur $\mathbb Z$; il est de plus \uline{fini} au dessus de $\mathbb Z[\frac1{\delta!}]$ -- i.e. aux voisinages en car. $p>\delta$ ($\delta=$ complexité de $\mathcal C_0$)
